%% ch14 --- trace-assert, complete: the same instrument in three languages
%%
%% Stub, generated by tools/gen_stubs.py from tools/chapters.json.
%% The brief below is the contract for this chapter: what it must argue, what
%% it must buy the reader, and what it must NOT repeat from the companion
%% volumes. Writing the chapter means deleting the stub block below (the one
%% that opens on the next \chapter line) and replacing it with the chapter
%% text -- after which this file is never regenerated.
%% If the brief turns out to be wrong once the code has been run, change it in
%% the manifest and record the contradiction in CLAUDE.md under *Resolved
%% questions*.

\chapter{trace-assert, complete: the same instrument in three languages}
\label{ch:traceassert}

\chapterstub{%
Argument: a deterministic assertion over an execution trace is the cheapest
evaluation an agent can have, and the reader can now build it in Python as
well as they could in C\#. The assertion types of agent-eval-bench's first
layer, ported to pytest \dash{} and the list is that project's own: it must
be copied from its specification when this chapter is written, never
reconstructed from memory, and the count printed here is whatever that
specification says. Each assertion as a fixture-backed function whose
failure message names the offending event; trace capture from the chapter 13
client; the comparison table across .NET, TypeScript and Python \dash{} what
each port needed in types, async, test framework and packaging; and
publishing \code{trace-assert} with \code{uv build} and trusted publishing
as the book's last exercise. Payoff: the artefact is a package, a blog post
and an interview answer at once. No new measurement; the suite's own run
time is reported. Three exercises. Every listing here is the real package
under \code{code/src/trace\_assert}, pulled in by region.

\medskip
\textbf{Word budget:} 3000.\\
\textbf{Ships in:} v1.0.\\
\textbf{Guiding project:} trace-assert final: the full assertion set, packaged and published.
}
